\documentclass[a4paper,11pt]{article}

\usepackage{polyglossia}
\setmainlanguage{greek}
\setotherlanguage{english}
\usepackage{fontspec}
\newfontfamily\greekfonttt{FreeMono}
\setmainfont{Linux Libertine O}
\usepackage{amsmath,amssymb}
\usepackage{geometry}
\geometry{margin=2.5cm}

\begin{document}

\title{Σημειώσεις Θεωρίας\\[4pt]
       Αναγνώριση Προτύπων}
\author{Μιχαήλ-Αθανάσιος Πέππας}
\maketitle

\tableofcontents
\clearpage

% ======================================================
\section{Μέρος 1: Ταξινόμηση στη Μία Διάσταση}
% ======================================================

% ------------------------------------------------------
\subsection{Εισαγωγή: πρόβλημα, οικογένειες μοντέλων και βασικές έννοιες}
% ------------------------------------------------------

Μελετάμε ένα κλασικό πρόβλημα \textbf{δυαδικής ταξινόμησης} με ένα μόνο χαρακτηριστικό:
\[
x \in \mathbb{R}, \quad y \in \{A,B\}.
\]

Το σενάριο αυτό είναι σκόπιμα απλό, αλλά θεωρητικά πολύ πλούσιο, γιατί επιτρέπει να συγκρίνουμε \textbf{διαφορετικές οικογένειες μοντέλων ταξινόμησης}:

\begin{itemize}
  \item \textbf{Μη παραμετρικές} μέθοδοι, όπως ο $k$-Nearest Neighbors ($k$-NN).
  \item \textbf{Γραμμικοί διακριτικοί ταξινομητές} (Perceptron, LMS).
  \item \textbf{Γενετικά μοντέλα} με Γκαουσιανές ανά κλάση και ταξινομητή Bayes.
\end{itemize}

Στόχος μας είναι \textbf{θεωρητικός}: να καταλάβουμε
\begin{itemize}
  \item πώς διαφέρουν αυτές οι οικογένειες,
  \item τι υποθέσεις κάνουν για τα δεδομένα,
  \item πότε προσεγγίζουν τον \textbf{βέλτιστο ταξινομητή Bayes}.
\end{itemize}

Βασικές έννοιες:

\begin{description}
  \item[Discriminative vs generative]
    \begin{itemize}
      \item \textbf{Generative} μοντέλα: μοντελοποιούν \emph{ολόκληρη} την κατανομή
            \[
              p(x,y) = P(C)\,p(x\mid C),
            \]
            π.χ.\ Γκαουσιανές ανά κλάση, HMM. Από αυτά
            προκύπτει ο ταξινομητής μέσω του κανόνα Bayes.
      \item \textbf{Discriminative} μοντέλα: μοντελοποιούν άμεσα την απόφαση
            ή την $p(y\mid x)$, χωρίς να ορίζουν ρητά $p(x\mid y)$.
            Παραδείγματα: Perceptron, LMS, Logistic Regression, Νευρωνικά Δίκτυα.
    \end{itemize}

  \item[Parametric vs non-parametric]
    \begin{itemize}
      \item \textbf{Parametric}: η ``γνώση'' συμπυκνώνεται σε έναν μικρό
            αριθμό παραμέτρων (π.χ.\ μέσοι, διασπορές, βάρη).
      \item \textbf{Non-parametric}: δεν υπάρχει προκαθορισμένος περιορισμένος
            αριθμός παραμέτρων∙ το μοντέλο ``μεγαλώνει'' με τα δεδομένα (π.χ.\ $k$-NN).
    \end{itemize}

  \item[Γραμμικά vs μη γραμμικά όρια απόφασης]
    \begin{itemize}
      \item Στη 1-D, \textbf{γραμμικό όριο} σημαίνει ένα μόνο σημείο κατωφλίου $x^\ast$.
      \item Μη γραμμικό όριο μπορεί να αποτελείται από πολλά διαστήματα
            (π.χ.\ $A,B,A,B,\dots$), όπως στον $1$-NN.
    \end{itemize}

  \item[Bayes βέλτιστος ταξινομητής]
    \begin{itemize}
      \item Θεωρητικά, ο καλύτερος ταξινομητής είναι αυτός που
            \emph{ελαχιστοποιεί την πιθανότητα λάθους}:
            \[
              y^\ast(x) = \arg\max_{C} P(C\mid x).
            \]
      \item Το αντίστοιχο ελάχιστο σφάλμα ονομάζεται \textbf{Bayes σφάλμα}.
    \end{itemize}
\end{description}

Η ταξινόμηση στη μία διάσταση είναι έτσι μια ``μικρογραφία'' όλων αυτών των ιδεών:
κάθε μέθοδος δίνει διαφορετικό σχήμα για το όριο απόφασης και κάνει
διαφορετικές υποθέσεις για το πώς προκύπτουν τα δεδομένα.

% ------------------------------------------------------
\subsection{Μη παραμετρικοί ταξινομητές: $k$-Nearest Neighbors}
% ------------------------------------------------------

\subsubsection*{Διαισθητική ιδέα και χαρακτηριστικά}

Ο ταξινομητής $k$-Nearest Neighbors ($k$-NN):

\begin{itemize}
  \item \textbf{Δεν μαθαίνει ρητές παραμέτρους} (όπως βάρη ή μέσους).
  \item Αποθηκεύει όλα τα εκπαιδευτικά ζεύγη
        $\{(x_n,y_n)\}_{n=1}^N$.
  \item Για ένα νέο $x$, βρίσκει τους $k$ κοντινότερους γείτονες
        (στη 1-D: με βάση την απόσταση $|x - x_n|$).
  \item \textbf{Απόφαση}: επιλέγει την κλάση που εμφανίζεται πιο συχνά
        μεταξύ αυτών των $k$ δειγμάτων.
\end{itemize}

Άρα:
\begin{itemize}
  \item Είναι \textbf{μη παραμετρικός}: δεν υποθέτει συγκεκριμένη μορφή για $p(x\mid C)$.
  \item Είναι \textbf{βασισμένος σε παραδείγματα} (instance-based):
        η γνώση μένει στα ίδια τα παραδείγματα, όχι σε ένα μικρό διάνυσμα παραμέτρων.
\end{itemize}

\subsubsection*{Όρια απόφασης στη μία διάσταση}

Στη 1-D τα δεδομένα τοποθετούνται πάνω στον άξονα των πραγματικών.

\begin{itemize}
  \item Για $1$-NN:
        \begin{itemize}
          \item Κάθε σημείο $x$ ταξινομείται στην κλάση του \emph{πλησιέστερου} δείγματος.
          \item Τα \textbf{όρια απόφασης} είναι τα \emph{midpoints} μεταξύ
                διαδοχικών δειγμάτων διαφορετικής κλάσης.
          \item Το όριο απόφασης είναι πολύ ``οδοντωτό'' και προσαρμόζεται ακριβώς
                σε κάθε σημείο του training set.
        \end{itemize}
  \item Για $k>1$ (π.χ.\ $3$-NN):
        \begin{itemize}
          \item Η απόφαση είναι με \textbf{πλειοψηφία}.
          \item Τα όρια απόφασης γίνονται πιο ``ομαλά'' (λιγότερες εναλλαγές κλάσης).
        \end{itemize}
\end{itemize}

Από την οπτική bias–variance:
\begin{itemize}
  \item \textbf{Μικρό $k$}: χαμηλό bias, υψηλό variance (έντονη προσαρμογή και σε θόρυβο).
  \item \textbf{Μεγάλο $k$}: υψηλότερο bias, χαμηλότερο variance (πιο ``γενική'' εικόνα).
\end{itemize}

\subsubsection*{Θεωρητική συμπεριφορά και συγκρίσεις}

\begin{itemize}
  \item Θεωρητικά, για $N\to\infty$ και κατάλληλη επιλογή $k$ (που αυξάνει με $N$),
        ο $k$-NN μπορεί να προσεγγίσει τον Bayes ταξινομητή.
  \item Υπολογιστικά, είναι \textbf{ακριβός στο test}: για κάθε νέο $x$
        χρειάζεται να βρει τις αποστάσεις από όλα τα εκπαιδευτικά δείγματα.
  \item Σε αντίθεση με τα γραμμικά μοντέλα (Perceptron/LMS),
        \textbf{μπορεί να παράγει μη γραμμικά όρια απόφασης} ακόμη και στη 1-D
        (πολλά διαστήματα).
  \item Δεν δίνει άμεσα θεωρητική ερμηνεία μέσω παραμέτρων (π.χ.\ δεν μαθαίνει μέσους/διασπορές).
\end{itemize}

% ------------------------------------------------------
\subsection{Γραμμικοί διακριτικοί ταξινομητές: Perceptron και LMS}
% ------------------------------------------------------

Και τα δύο μοντέλα ανήκουν στην ίδια οικογένεια: \textbf{γραμμικοί διακριτικοί ταξινομητές}.
Στη μία διάσταση, το μοντέλο έχει τη μορφή
\[
y(x) = w x + b.
\]
Στο τέλος, η απόφαση προκύπτει συγκρίνοντας το $y(x)$ με κάποιο κατώφλι.

\subsubsection*{Perceptron}

\paragraph{Μοντέλο και όριο απόφασης.}

Στην απλούστερη μορφή:
\[
\hat{y}(x) = \text{sign}(w x + b),
\]
όπου $\hat{y}\in\{-1,+1\}$ κωδικοποιεί τις δύο κλάσεις ($A,B$).
Το \textbf{όριο απόφασης} είναι
\[
w x + b = 0 \;\;\Rightarrow\;\; x^\ast = -\frac{b}{w},
\]
ένα \emph{μοναδικό} σημείο διαχωρισμού στη 1-D (γραμμικό όριο).

\paragraph{Κανόνας εκμάθησης (Perceptron rule).}

Για κάθε δείγμα $(x_n,t_n)$ με $t_n\in\{-1,+1\}$:
\begin{itemize}
  \item Υπολογίζουμε $\hat{y}_n = \text{sign}(w x_n + b)$.
  \item Αν $\hat{y}_n = t_n$, δεν κάνουμε τίποτα.
  \item Αν $\hat{y}_n \neq t_n$, ενημερώνουμε:
        \[
          w \leftarrow w + \eta\, t_n x_n, \qquad
          b \leftarrow b + \eta\, t_n,
        \]
        όπου $\eta>0$ είναι ο ρυθμός μάθησης.
\end{itemize}

Μια \textbf{epoch} = ένα πλήρες πέρασμα από όλα τα δείγματα.

\paragraph{Θεωρία και ιδιότητες.}

\begin{itemize}
  \item \textbf{Perceptron convergence theorem}: αν τα δεδομένα είναι
        \textbf{γραμμικά διαχωρίσιμα}, το perceptron συγκλίνει σε πεπερασμένα βήματα
        σε κάποιο $(w,b)$ που διαχωρίζει τέλεια τις δύο κλάσεις.
  \item Αν τα δεδομένα \textbf{δεν} είναι διαχωρίσιμα (π.χ.\ αλληλοεπικαλυπτόμενες Γκαουσιανές),
        το perceptron δεν συγκλίνει∙ συνεχίζει να τροποποιεί τα βάρη.
  \item Είναι \textbf{διακριτικό} και \textbf{παραμετρικό} μοντέλο:
        μαθαίνει απευθείας ένα γραμμικό όριο, χωρίς να μοντελοποιεί $p(x\mid C)$.
\end{itemize}

\subsubsection*{LMS (Widrow–Hoff): γραμμική παλινδρόμηση για ταξινόμηση}

Ο LMS αντιμετωπίζει την ταξινόμηση ως \textbf{γραμμική παλινδρόμηση}:
\[
y(x) = w x + b,
\]
και επιλέγει τιμές στόχου $t_n$ (π.χ.\ $t_n\in\{-1,+1\}$) έτσι ώστε
$y(x_n)\approx t_n$. Το όριο απόφασης προκύπτει πάλι από το σημείο όπου
$y(x)=0$ ή από ένα κατάλληλο κατώφλι.

\paragraph{Συνάρτηση κόστους.}

Για εκπαίδευση χρησιμοποιείται η \textbf{μέση τετραγωνική απόκλιση} (MSE):
\[
J(w,b) = \frac{1}{2}\sum_{n=1}^N (t_n - (w x_n + b))^2.
\]

\paragraph{Gradient descent και κανόνας ενημέρωσης.}

Παίρνοντας παραγώγους:
\[
\frac{\partial J}{\partial w}
= -\sum_{n=1}^N (t_n - y_n) x_n,\qquad
\frac{\partial J}{\partial b}
= -\sum_{n=1}^N (t_n - y_n),
\]
όπου $y_n = w x_n + b$.

Με \textbf{online} gradient descent (ένα δείγμα τη φορά):
\[
w \leftarrow w + \eta (t_n - y_n) x_n, \qquad
b \leftarrow b + \eta (t_n - y_n).
\]

\paragraph{Συγκρίσεις Perceptron vs LMS.}

\begin{itemize}
  \item Και τα δύο δίνουν \textbf{γραμμικό όριο απόφασης} στη 1-D.
  \item Perceptron: ενημερώνει μόνο αν υπάρχει \emph{λάθος ταξινόμηση},
        δεν ενδιαφέρεται για το ``πόσο'' λάθος είναι.
  \item LMS: ελαχιστοποιεί το \textbf{μέγεθος του σφάλματος} (MSE),
        άρα είναι πιο ``ομαλό'' και σταθερό σε θορυβώδη δεδομένα.
  \item Και τα δύο είναι \textbf{διακριτικά, παραμετρικά} μοντέλα και
        δουλεύουν καλά όταν το πραγματικό όριο απόφασης είναι περίπου γραμμικό.
\end{itemize}

% ------------------------------------------------------
\subsection{Γενετικά Γκαουσιανά μοντέλα και Ταξινομητής Bayes}
% ------------------------------------------------------

\subsubsection*{Γενετικό μοντέλο με Γκαουσιανές ανά κλάση}

Υποθέτουμε ότι για κάθε κλάση $C\in\{A,B\}$ η κατανομή του $x$ είναι
Γκαουσιανή:
\[
p(x \mid C) = \mathcal{N}(x \mid \mu_C,\sigma_C^2)
  = \frac{1}{\sqrt{2\pi\sigma_C^2}}
    \exp\biggl( -\frac{(x-\mu_C)^2}{2\sigma_C^2} \biggr).
\]

Επιπλέον, έχουμε priors $P(A)$, $P(B)$ με $P(A)+P(B)=1$.
Το μοντέλο είναι \textbf{γενετικό}: περιγράφει πώς ``παράγονται'' τα δεδομένα.

\subsubsection*{ML εκτίμηση για μέσο και διασπορά}

Για τα δείγματα της κλάσης $C$,
$\{x_1,\dots,x_{N_C}\}$, η \textbf{μέγιστη πιθανοφάνεια} (ML) δίνει:

\[
\mu_C^{\text{ML}} = \frac{1}{N_C}\sum_{n=1}^{N_C} x_n,
\]
\[
\sigma_C^{2\,\text{ML}} = \frac{1}{N_C}\sum_{n=1}^{N_C}(x_n - \mu_C^{\text{ML}})^2.
\]

Πρόκειται για τις αναμενόμενες εκφράσεις: ο δειγματικός μέσος και η
δειγματική διασπορά με παρονομαστή $N_C$ (ML, όχι unbiased).

\subsubsection*{Ταξινομητής Bayes και όριο απόφασης}

Ο ταξινομητής Bayes αποφασίζει υπέρ της κλάσης με τη μεγαλύτερη posterior:
\[
\text{απόφαση}(x)=
\begin{cases}
A, & \text{αν }P(A\mid x) > P(B\mid x),\\
B, & \text{αλλιώς.}
\end{cases}
\]

Με τον κανόνα Bayes:
\[
P(A\mid x) = \frac{P(A)p(x\mid A)}{p(x)},\quad
P(B\mid x) = \frac{P(B)p(x\mid B)}{p(x)},
\]
όπου $p(x)=P(A)p(x\mid A)+P(B)p(x\mid B)$.

Το \textbf{όριο απόφασης} $x^\ast$ ικανοποιεί:
\[
P(A)p(x^\ast\mid A) = P(B)p(x^\ast\mid B).
\]

Αν οι δύο Γκαουσιανές έχουν \emph{ίση διασπορά} $\sigma_A^2=\sigma_B^2=\sigma^2$,
τότε μετά από λογαρίθμηση παίρνουμε \emph{γραμμικό} όρο στο $x$ και άρα
ένα μοναδικό σημείο:
\[
x^\ast = \frac{2\sigma^2 \ln\!\Bigl(\dfrac{P(B)}{P(A)}\Bigr)
    - (\mu_A^2 - \mu_B^2)}{2(\mu_B - \mu_A)}.
\]

Ειδική περίπτωση $P(A)=P(B)$:
\[
\boxed{x^\ast = \frac{\mu_A+\mu_B}{2}.}
\]

\subsubsection*{Bayes σφάλμα και απόσταση των Γκαουσιανών}

Ο Bayes ταξινομητής ελαχιστοποιεί την πιθανότητα σφάλματος
\[
P_e = \int_{\mathcal{R}_A} P(B)p(x\mid B)\,dx
    + \int_{\mathcal{R}_B} P(A)p(x\mid A)\,dx,
\]
ή ισοδύναμα
\[
P_e = \int_{-\infty}^{+\infty} \min\{P(A)p(x\mid A),P(B)p(x\mid B)\}\,dx.
\]

Για δύο Γκαουσιανές με ίσες διασπορές $\sigma^2$ και ίσους priors,
αν ορίσουμε
\[
d = \frac{|\mu_B-\mu_A|}{\sigma},
\]
τότε (χωρίς απόδειξη):
\[
\boxed{P_e = \Phi\!\Bigl(-\frac{d}{2}\Bigr)},
\]
όπου $\Phi$ η CDF της $\mathcal{N}(0,1)$. Όσο μεγαλύτερο $d$, τόσο πιο
καλά διαχωρίζονται οι κλάσεις και τόσο μικρότερο το Bayes σφάλμα.

\subsubsection*{Συγκρίσεις με διακριτικούς ταξινομητές}

\begin{itemize}
  \item Τα γενετικά Γκαουσιανά μοντέλα \textbf{μοντελοποιούν το $p(x\mid C)$},
        κάτι που επιτρέπει όχι μόνο ταξινόμηση αλλά και ``μοντελοποίηση των δεδομένων''.
  \item Αν οι πραγματικές κατανομές είναι κοντά σε Γκαουσιανές,
        το μοντέλο μπορεί να πλησιάσει πολύ τον βέλτιστο Bayes ταξινομητή.
  \item Σε αντίθεση, Perceptron και LMS δεν έχουν έννοια ``κατανομής'',
        αλλά μαθαίνουν \emph{άμεσα} το όριο απόφασης.
  \item Στη 1-D, τόσο το Γκαουσιανό μοντέλο (με ίσες διασπορές) όσο και ο
        γραμμικός ταξινομητής παράγουν ουσιαστικά ένα \emph{μοναδικό κατώφλι}.
\end{itemize}

% ------------------------------------------------------
\subsection{Συγκεντρωτική σύγκριση μεθόδων (Μέρος 1)}
% ------------------------------------------------------

\begin{center}
\begin{tabular}{l|c|c|c|c}
\textbf{Μέθοδος} & \textbf{Generative} & \textbf{Parametric} & \textbf{Όριο στη 1-D} & \textbf{Ευελιξία} \\
\hline
$k$-NN              & όχι   & όχι   & πολλαπλά διαστήματα & πολύ υψηλή (τοπικό) \\
Perceptron          & όχι   & ναι   & ένα κατώφλι         & μέτρια, γραμμικό \\
LMS                 & όχι   & ναι   & ένα κατώφλι         & μέτρια, γραμμικό \\
Γκαουσιανό + Bayes  & ναι   & ναι   & συνήθως ένα κατώφλι & υψηλή αν υπόθεση σωστή \\
\end{tabular}
\end{center}

Σημεία-κλειδιά:
\begin{itemize}
  \item $k$-NN: δυνατός μη παραμετρικός ταξινομητής, αλλά ακριβός στο test
        και δύσκολος σε μεγάλες διαστάσεις.
  \item Perceptron/LMS: απλά διακριτικά, παραμετρικά μοντέλα∙
        χρήσιμα για κατανόηση γραμμικών ορίων και gradient descent.
  \item Γκαουσιανό generative μοντέλο: συνδέει ταξινόμηση με
        πιθανοτικές υποθέσεις και Bayes βέλτιστοτητα.
\end{itemize}

\clearpage

% ======================================================
\section{Μέρος 2: Κρυφά Μοντέλα Markov (HMM)}
% ======================================================

% ------------------------------------------------------
\subsection{Θεωρητικό πλαίσιο: στοχαστικές ακολουθίες και ιδιότητα Markov}
% ------------------------------------------------------

Σε πολλά προβλήματα (ομιλία, βιολογικές ακολουθίες, NLP) η \textbf{τάξη}
των παρατηρήσεων παίζει κεντρικό ρόλο:
\[
O_0,O_1,\dots,O_T.
\]
Είναι ανεπαρκές να αντιμετωπίζουμε τα $O_t$ ως ανεξάρτητα∙ χρειάζεται
μοντέλο που να περιγράφει πώς εξελίσσεται η κατάσταση στο χρόνο.

Η \textbf{ιδιότητα Markov πρώτης τάξης} λέει ότι:
\[
P(q_t \mid q_{0:t-1}) = P(q_t \mid q_{t-1}),
\]
δηλαδή η επόμενη κατάσταση εξαρτάται μόνο από την \emph{τρέχουσα}.

Τα HMM συνδυάζουν:
\begin{itemize}
  \item \textbf{αλυσιδωτή δομή} Markov για τις κρυφές καταστάσεις,
  \item \textbf{πιθανότητες εκπομπής} για τις παρατηρήσεις.
\end{itemize}

% ------------------------------------------------------
\subsection{Ορισμός HMM ως γενετικού μοντέλου}
% ------------------------------------------------------

Ένα HMM ορίζεται από τις παραμέτρους $\lambda = (\pi, A, B)$:

\begin{itemize}
  \item \textbf{Καταστάσεις} $q_t \in \{1,\dots,S\}$.
  \item \textbf{Παρατηρήσεις} $O_t$ (διακριτές ή συνεχείς).
  \item \textbf{Αρχική κατανομή} $\pi_i = P(q_0=i)$.
  \item \textbf{Μεταβατικές πιθανότητες} $a_{ij} = P(q_t=j \mid q_{t-1}=i)$.
  \item \textbf{Πιθανότητες εκπομπής} $b_j(o) = P(O_t=o \mid q_t=j)$.
\end{itemize}

\paragraph{Γενετική διαδικασία.}

\begin{enumerate}
  \item Επιλέγουμε αρχική κατάσταση $q_0$ από την $\pi$.
  \item Για $t=0,1,\dots,T$:
        \begin{itemize}
          \item εκπέμπουμε παρατήρηση $O_t$ από την κατανομή $b_{q_t}(\cdot)$,
          \item μεταβαίνουμε σε νέα κατάσταση $q_{t+1}$ με πιθανότητα $a_{q_t,q_{t+1}}$.
        \end{itemize}
\end{enumerate}

Η \textbf{joint} πιθανότητα διαδρομής και παρατηρήσεων είναι:
\[
P(q_0,\dots,q_T, O_0,\dots,O_T)
= \pi_{q_0}\, b_{q_0}(O_0)\, \prod_{t=1}^T a_{q_{t-1},q_t}\, b_{q_t}(O_t).
\]

Όπως τα Γκαουσιανά μοντέλα στο Μέρος~1, έτσι κι εδώ έχουμε ένα \textbf{γενετικό μοντέλο},
μόνο που τώρα το $x$ αντικαθίσταται από ολόκληρη την ακολουθία $O_{0:T}$.

% ------------------------------------------------------
\subsection{Τα τρία βασικά προβλήματα σε ένα HMM}
% ------------------------------------------------------

Για ένα HMM και μια ακολουθία παρατηρήσεων $O_{0:T}$, τα κλασικά θεωρητικά
προβλήματα είναι:

\begin{enumerate}
  \item \textbf{Πρόβλημα αξιολόγησης}:
        \[
        \text{να υπολογίσουμε } P(O_{0:T} \mid \lambda).
        \]
  \item \textbf{Πρόβλημα αποκωδικοποίησης}:
        \[
        \text{να βρούμε την πιο πιθανή ακολουθία καταστάσεων } q_{0:T}^\ast.
        \]
  \item \textbf{Πρόβλημα εκμάθησης}:
        \[
        \text{να εκτιμήσουμε τις παραμέτρους } \lambda
        \text{ από παρατηρούμενες ακολουθίες}.
        \]
\end{enumerate}

Τα αντίστοιχα βασικά εργαλεία είναι:
\begin{itemize}
  \item \textbf{Forward} (ή forward–backward) αλγόριθμος,
  \item \textbf{Viterbi} αλγόριθμος,
  \item \textbf{EM / Baum–Welch} (και παραλλαγές όπως Viterbi-training).
\end{itemize}

% ------------------------------------------------------
\subsection{Πρόβλημα αξιολόγησης: forward αλγόριθμος}
% ------------------------------------------------------

Η ``άμεση'' λύση του $P(O_{0:T})$ είναι:
\[
P(O_{0:T}) = \sum_{q_{0:T}} P(q_{0:T},O_{0:T}),
\]
άθροισμα σε όλες τις δυνατές διαδρομές καταστάσεων∙ αυτό όμως είναι
\textbf{εκθετικό} σε $T$.

Ο \textbf{forward αλγόριθμος} χρησιμοποιεί \textbf{δυναμικό προγραμματισμό}.

Ορίζουμε:
\[
\alpha_t(j) = P(O_0,\dots,O_t, q_t=j).
\]

\paragraph{Αρχικοποίηση.}
\[
\alpha_0(j) = \pi_j\, b_j(O_0), \quad j=1,\dots,S.
\]

\paragraph{Επανάληψη.}
Για $t=1,\dots,T$:
\[
\alpha_t(j) = b_j(O_t) \sum_{i=1}^S \alpha_{t-1}(i)\, a_{ij}.
\]

\paragraph{Τερματισμός.}
\[
P(O_{0:T}) = \sum_{j=1}^S \alpha_T(j).
\]

Ερμηνεία:
\begin{itemize}
  \item Η $\alpha_t(j)$ είναι η πιθανότητα να έχουμε παρατηρήσει
        $O_0,\dots,O_t$ και να βρισκόμαστε στην κατάσταση $j$ στο $t$.
  \item Ο αλγόριθμος ``συγκεντρώνει'' όλες τις διαδρομές που καταλήγουν σε κάθε $j$
        χωρίς να τις απαριθμεί μία-μία.
\end{itemize}

% ------------------------------------------------------
\subsection{Πρόβλημα αποκωδικοποίησης: αλγόριθμος Viterbi}
% ------------------------------------------------------

Τώρα θέλουμε την \textbf{πιο πιθανή διαδρομή καταστάσεων}:
\[
q_0^\ast,\dots,q_T^\ast = \arg\max_{q_{0:T}} P(q_{0:T},O_{0:T}).
\]

Ο \textbf{Viterbi} αλγόριθμος μοιάζει με τον forward, αλλά αντί για
άθροισμα κρατάει το \emph{μέγιστο}.

Ορίζουμε:
\[
\delta_t(j) = \max_{q_{0:t-1}}
P(q_{0:t-1}, q_t=j, O_{0:t}),
\]
και δείκτες backtracking
\[
\psi_t(j) = \arg\max_i \left[ \delta_{t-1}(i) a_{ij} \right].
\]

\paragraph{Αρχικοποίηση.}
\[
\delta_0(j) = \pi_j b_j(O_0),\quad
\psi_0(j) = 0.
\]

\paragraph{Επανάληψη.}
\[
\delta_t(j) = b_j(O_t)\, \max_i [\delta_{t-1}(i) a_{ij}],
\quad
\psi_t(j) = \arg\max_i [\delta_{t-1}(i) a_{ij}].
\]

\paragraph{Τερματισμός \& backtracking.}
\[
P^\ast = \max_j \delta_T(j),\quad
q_T^\ast = \arg\max_j \delta_T(j),
\]
και για $t=T-1,\dots,0$:
\[
q_t^\ast = \psi_{t+1}(q_{t+1}^\ast).
\]

Σημαντική διαφορά:
\begin{itemize}
  \item Forward: υπολογίζει \textbf{ολική πιθανότητα} $P(O_{0:T})$ (άθροισμα όλων των διαδρομών).
  \item Viterbi: υπολογίζει \textbf{μια διαδρομή μέγιστης πιθανότητας} (MAP path).
\end{itemize}

% ------------------------------------------------------
\subsection{Πρόβλημα εκμάθησης: EM / Baum–Welch}
% ------------------------------------------------------

Θέλουμε να εκτιμήσουμε $\lambda = (\pi, A, B)$ από παρατηρήσεις
$\mathcal{D} = \{O^{(m)}_{0:T_m}\}_m$.

Οι καταστάσεις $q_t$ είναι κρυφές, άρα η άμεση ML βελτιστοποίηση είναι δύσκολη.
Χρησιμοποιούμε τον \textbf{αλγόριθμο EM} στην ειδική μορφή του για HMM,
γνωστό ως \textbf{Baum–Welch}.

\subsubsection*{E-step: forward–backward, $\gamma$ και $\xi$}

Εκτός από τα forward $\alpha_t(j)$, ορίζουμε τα backward:
\[
\beta_t(j) = P(O_{t+1},\dots,O_T \mid q_t=j).
\]

\paragraph{Αρχικοποίηση.}
\[
\beta_T(j) = 1.
\]

\paragraph{Recursion.}
\[
\beta_t(i) = \sum_{j=1}^S a_{ij} b_j(O_{t+1}) \beta_{t+1}(j).
\]

Με αυτά, υπολογίζουμε:
\[
\gamma_t(i) = P(q_t=i\mid O_{0:T})
= \frac{\alpha_t(i)\beta_t(i)}{P(O_{0:T})},
\]
\[
\xi_t(i,j) = P(q_t=i, q_{t+1}=j \mid O_{0:T})
= \frac{\alpha_t(i) a_{ij} b_j(O_{t+1}) \beta_{t+1}(j)}{P(O_{0:T})}.
\]

Ερμηνείες:
\begin{itemize}
  \item $\gamma_t(i)$: αναμενόμενη ``παρουσία'' στην κατάσταση $i$ στο χρόνο $t$.
  \item $\xi_t(i,j)$: αναμενόμενος αριθμός μεταβάσεων $i\to j$ μεταξύ $t$ και $t+1$.
\end{itemize}

\subsubsection*{M-step: ενημέρωση παραμέτρων}

\begin{itemize}
  \item Αρχικές πιθανότητες:
        \[
        \pi_i^{\text{new}} = \gamma_0(i).
        \]
  \item Μεταβατικές πιθανότητες:
        \[
        a_{ij}^{\text{new}}
        = \frac{\sum_{t=0}^{T-1} \xi_t(i,j)}{\sum_{t=0}^{T-1} \gamma_t(i)}.
        \]
  \item Πιθανότητες εκπομπής (διακριτές παρατηρήσεις):
        \[
        b_j^{\text{new}}(o)
        = \frac{\sum_{t:O_t=o} \gamma_t(j)}{\sum_{t=0}^T \gamma_t(j)}.
        \]
\end{itemize}

Το EM αυξάνει μονοτονικά το log-likelihood, αλλά μπορεί να καταλήξει
σε \textbf{τοπικά} optima.

% ------------------------------------------------------
\subsection{Viterbi-training (hard EM) και σύγκριση}
% ------------------------------------------------------

Το Viterbi-training είναι μια απλοποιημένη (hard) εκδοχή του EM:

\begin{enumerate}
  \item \textbf{Hard E-step}: τρέχουμε Viterbi για να βρούμε μια διαδρομή
        $q_{0:T}^\ast$ μέγιστης πιθανότητας.
  \item \textbf{M-step}: θεωρούμε αυτή τη διαδρομή ως ``παρατηρημένη'':
        μετράμε
        \begin{itemize}
          \item πόσες φορές βρισκόμαστε σε κάθε κατάσταση,
          \item πόσες μεταβάσεις $i\to j$ γίνονται,
          \item από κάθε κατάσταση πόσες φορές εκπέμπεται κάθε σύμβολο,
        \end{itemize}
        και ενημερώνουμε $A,B$ ως σχετικές συχνότητες.
\end{enumerate}

Σύγκριση:
\begin{itemize}
  \item \textbf{Baum–Welch (forward–backward)}:
        \begin{itemize}
          \item soft counts (χρησιμοποιεί όλες τις διαδρομές, σταθμισμένες με τις πιθανότητές τους),
          \item στατιστικά πιο ορθό (πραγματικό EM),
          \item καλύτερη συμπεριφορά όταν το μοντέλο είναι ασαφές/οι διαδρομές πολλές.
        \end{itemize}
  \item \textbf{Viterbi-training}:
        \begin{itemize}
          \item hard assignment: χρησιμοποιεί μόνο την καλύτερη διαδρομή,
          \item υπολογιστικά πιο απλό και γρήγορο,
          \item μπορεί να παγιδευτεί πιο εύκολα σε ``κακές'' λύσεις.
        \end{itemize}
\end{itemize}

% ------------------------------------------------------
\subsection{Σχέση HMM με τα μοντέλα του Μέρους 1}
% ------------------------------------------------------

\begin{itemize}
  \item Όπως οι Γκαουσιανές στο Μέρος~1, τα HMM είναι \textbf{γενετικά μοντέλα}:
        δίνουν ένα πλήρες $p(O_{0:T},q_{0:T})$.
  \item Όπως στο 1-D μοντέλο μάθαμε παραμέτρους με ML (μέσους, διασπορές),
        εδώ μαθαίνουμε παραμέτρους $\lambda$ με EM.
  \item Η λογική ``\textbf{latent μεταβλητές + EM}'' είναι κοινή:
        \begin{itemize}
          \item στο 1-D μπορούμε να έχουμε άγνωστες κλάσεις,
          \item στο HMM έχουμε κρυφές καταστάσεις σε κάθε χρονική στιγμή.
        \end{itemize}
  \item Σε αντίθεση με Perceptron/LMS, τα HMM δεν δίνουν άμεσα ένα
        \emph{decision boundary} σε χώρο χαρακτηριστικών, αλλά ένα
        \emph{πιθανοτικό μοντέλο ακολουθιών}, από το οποίο μπορούμε να
        παράγουμε ταξινομητές (π.χ.\ μοντέλα ανά λέξη στην αναγνώριση ομιλίας).
\end{itemize}

\clearpage

% ======================================================
\section{Μέρος 3: Backpropagation σε Απλό Νευρωνικό Δίκτυο}
% ======================================================

Σκοπός αυτού του μέρους είναι να κατανοήσουμε σε θεωρητικό επίπεδο
το \textbf{backpropagation}:

\begin{itemize}
  \item πώς συνδέεται με τον \textbf{κανόνα αλυσίδας} στις παραγώγους,
  \item πώς γενικεύει την ιδέα του \textbf{gradient descent} από γραμμικά
        μοντέλα (LMS) σε \emph{πολυεπίπεδα} νευρωνικά δίκτυα,
  \item πώς τοπικές παράγωγοι συνδυάζονται για να δώσουν
        $\partial L / \partial w$ για κάθε βάρος.
\end{itemize}

% ------------------------------------------------------
\subsection{Νευρωνικά δίκτυα ως διακριτικά, μη γραμμικά μοντέλα}
% ------------------------------------------------------

Ένα απλό feedforward νευρωνικό δίκτυο (MLP) μπορεί να θεωρηθεί ως
σύνθεση:
\[
\text{είσοδος} \xrightarrow{\text{γραμμικός μετασχηματισμός}}
\xrightarrow{\text{μη γραμμική ενεργοποίηση}}
\xrightarrow{\text{γραμμικός μετασχηματισμός}} \dots \xrightarrow{} \hat{y}.
\]

Βασικά στοιχεία:
\begin{itemize}
  \item \textbf{Βάρη} και \textbf{bias} στα άκρα κάθε σύνδεσης.
  \item Μη γραμμική ενεργοποίηση (π.χ.\ logistic):
        \[
          \sigma(z) = \frac{1}{1+e^{-z}}.
        \]
  \item Συνάρτηση κόστους $L(y,\hat{y})$ (π.χ.\ MSE ή cross-entropy).
\end{itemize}

Σε αντίθεση με Perceptron/LMS:
\begin{itemize}
  \item Έχουμε \textbf{πολλαπλά επίπεδα} (hidden layers).
  \item Το μοντέλο μπορεί να προσεγγίσει \emph{πολύπλοκα μη γραμμικά}
        decision boundaries (θεωρία καθολικής προσέγγισης).
\end{itemize}

Παραμένει όμως \textbf{διακριτικό} μοντέλο: δίνει άμεσα $f(x)\approx P(y\mid x)$
και όχι $p(x\mid y)$.

% ------------------------------------------------------
\subsection{Υπολογιστικός γράφος και κανόνας αλυσίδας}
% ------------------------------------------------------

Το δίκτυο μπορεί να ιδωθεί ως \textbf{υπολογιστικός γράφος}:

\begin{itemize}
  \item Κάθε κόμβος υπολογίζει μια συνάρτηση των εισόδων του.
  \item Η έξοδος $\hat{y}$ είναι συνάρτηση όλων των βαρών.
  \item Το κόστος $L$ είναι τελικά συνάρτηση των βαρών:
        \[
        L = L(w_1,w_2,\dots).
        \]
\end{itemize}

Για να εκπαιδεύσουμε το δίκτυο με gradient descent, χρειαζόμαστε
\[
\frac{\partial L}{\partial w_i}
\quad \text{για όλα τα βάρη } w_i.
\]

Ο \textbf{κανόνας αλυσίδας} μας λέει ότι η παράγωγος ως προς ένα βάρος
ισούται με το γινόμενο όλων των κατάλληλων \emph{τοπικών παραγώγων}
που συνδέουν αυτό το βάρος με την έξοδο $L$.

Το backpropagation είναι απλώς ένας \textbf{αποδοτικός αλγόριθμος}
που:
\begin{itemize}
  \item κάνει ένα \textbf{forward pass} (υπολογίζει όλες τις ενδιάμεσες τιμές),
  \item και μετά ένα \textbf{backward pass} (υπολογίζει συστηματικά τις παραγώγους).
\end{itemize}

% ------------------------------------------------------
\subsection{Τυπικό 2-στρωσης δίκτυο και συνάρτηση κόστους}
% ------------------------------------------------------

Έστω δίκτυο με είσοδο $\mathbf{x}=(x_1,x_2,x_3,x_4)$, ένα κρυφό νευρώνα $h_1$
και μία έξοδο $\hat{y}$ (για απλότητα):

\begin{itemize}
  \item \textbf{Hidden layer}:
        \[
        s^{(h)}_1 = w_1 x_1 + w_2 x_2 + w_3 x_3 + w_4 x_4 + b^{(h)}_1,
        \]
        \[
        h_1 = \sigma(s^{(h)}_1).
        \]
  \item \textbf{Output layer} (συνοπτικά):
        \[
        s^{(y)} = v_1 h_1 + \dots + b^{(y)}, \qquad
        \hat{y} = \sigma(s^{(y)}).
        \]
\end{itemize}

Ως συνάρτηση κόστους παίρνουμε (για ένα δείγμα):
\[
L(y,\hat{y}) = (y - \hat{y})^2.
\]

Στόχος: να βρούμε
\[
\frac{\partial L}{\partial w_1},
\]
δηλαδή πόσο αλλάζει το κόστος αν μεταβάλουμε λίγο το βάρος $w_1$.

% ------------------------------------------------------
\subsection{Αναλυτικός υπολογισμός $\partial L / \partial w_1$}
% ------------------------------------------------------

Εφαρμόζουμε συστηματικά τον κανόνα αλυσίδας:

\[
\frac{\partial L}{\partial w_1}
= \frac{\partial L}{\partial \hat{y}}\,
  \frac{\partial \hat{y}}{\partial s^{(y)}}\,
  \frac{\partial s^{(y)}}{\partial h_1}\,
  \frac{\partial h_1}{\partial s^{(h)}_1}\,
  \frac{\partial s^{(h)}_1}{\partial w_1}.
\]

Υπολογίζουμε κάθε παράγοντα:

\begin{enumerate}
  \item Από $L = (y-\hat{y})^2$:
        \[
        \frac{\partial L}{\partial \hat{y}}
        = 2(\hat{y} - y).
        \]
  \item Logistic παράγωγος:
        \[
        \frac{\partial \hat{y}}{\partial s^{(y)}}
        = \sigma'(s^{(y)}) = \hat{y}(1-\hat{y}).
        \]
  \item Από $s^{(y)} = v_1 h_1 + \dots + b^{(y)}$:
        \[
        \frac{\partial s^{(y)}}{\partial h_1} = v_1.
        \]
  \item Από $h_1 = \sigma(s^{(h)}_1)$:
        \[
        \frac{\partial h_1}{\partial s^{(h)}_1}
        = \sigma'(s^{(h)}_1) = h_1(1-h_1).
        \]
  \item Από $s^{(h)}_1 = w_1 x_1 + \dots + b^{(h)}_1$:
        \[
        \frac{\partial s^{(h)}_1}{\partial w_1} = x_1.
        \]
\end{enumerate}

Συνδυάζοντας:
\[
\frac{\partial L}{\partial w_1}
= 2(\hat{y} - y)\,\hat{y}(1-\hat{y})\, v_1\,
  h_1(1-h_1)\, x_1.
\]

Αυτή είναι η \textbf{τυπική μορφή} του gradient:
\[
\text{gradient στο βάρος}
= (\text{σφάλμα στην έξοδο})
\cdot (\text{αλυσίδα τοπικών παραγώγων})
\cdot (\text{είσοδος στη σύνδεση}).
\]

% ------------------------------------------------------
\subsection{Διατύπωση με σήματα σφάλματος $\delta$ (error signals)}
% ------------------------------------------------------

Για ευκολία γράφουμε:

\begin{itemize}
  \item \textbf{Σφάλμα εξόδου}:
        \[
        \delta^{(y)}
        = \frac{\partial L}{\partial s^{(y)}}
        = \frac{\partial L}{\partial \hat{y}}\,
          \frac{\partial \hat{y}}{\partial s^{(y)}}
        = 2(\hat{y}-y)\,\hat{y}(1-\hat{y}).
        \]
  \item \textbf{Σφάλμα κρυφού νευρώνα}:
        \[
        \delta^{(h)}_1
        = \frac{\partial L}{\partial s^{(h)}_1}
        = \delta^{(y)}\, v_1\, h_1(1-h_1).
        \]
\end{itemize}

Τότε:
\[
\boxed{
\frac{\partial L}{\partial w_1}
= \delta^{(h)}_1\, x_1.
}
\]

Αυτό δείχνει πολύ καθαρά ότι:
\begin{itemize}
  \item τα \textbf{$\delta$ σφάλματα} διαδίδονται προς τα πίσω στο δίκτυο,
  \item κάθε gradient βάρους είναι απλώς \emph{σφάλμα στον κόμβο-στόχο}
        επί \emph{είσοδο στη συγκεκριμένη σύνδεση}.
\end{itemize}

% ------------------------------------------------------
\subsection{Συγκρίσεις με LMS, Perceptron και HMM}
% ------------------------------------------------------

\begin{itemize}
  \item Το backpropagation είναι ακριβώς η γενίκευση της ιδέας του
        \textbf{gradient descent} από γραμμικά μοντέλα (LMS) σε
        βαθιά, μη γραμμικά δίκτυα.
  \item Όπως στο LMS, ορίζουμε μια συνάρτηση κόστους $L$ και
        ενημερώνουμε βάρη $w \leftarrow w - \eta \frac{\partial L}{\partial w}$.
  \item Σε αντίθεση με το Perceptron, δεν βασιζόμαστε σε μονοτονική
        ενημέρωση μόνο όταν υπάρχει λάθος, αλλά σε \emph{συνεχή βελτιστοποίηση}
        του κόστους μέσω παραγώγων.
  \item Σε αντίθεση με τα γενετικά μοντέλα (Γκαουσιανές, HMM),
        τα MLP είναι καθαρά \textbf{διακριτικά}:
        δεν μοντελοποιούν $p(x\mid y)$, αλλά μάθουν έμμεσα $P(y\mid x)$.
\end{itemize}

% ------------------------------------------------------
\subsection{Σημεία-κλειδιά για θεωρητικές ερωτήσεις backprop}
% ------------------------------------------------------

\begin{itemize}
  \item Backpropagation = \textbf{κανόνας αλυσίδας} + \textbf{υπολογιστικός γράφος}.
  \item Απαιτεί:
        \begin{itemize}
          \item forward pass για τις ενδιάμεσες τιμές,
          \item backward pass για τα $\delta$ σφάλματα.
        \end{itemize}
  \item Για logistic:
        \[
        \sigma'(z) = \sigma(z)(1-\sigma(z)).
        \]
  \item Για MSE:
        \[
        \frac{\partial L}{\partial \hat{y}} = 2(\hat{y}-y).
        \]
  \item Κάθε gradient βάρους γράφεται ως:
        \[
        \frac{\partial L}{\partial w} = \delta_{\text{κόμβος προορισμού}} \cdot
        \text{είσοδος στη σύνδεση}.
        \]
  \item Η φιλοσοφία είναι ίδια με το LMS, αλλά με \emph{πολλά επίπεδα}
        και \emph{μη γραμμικές} ενεργοποιήσεις, άρα πολύ μεγαλύτερη εκφραστικότητα
        και ικανότητα προσέγγισης πολύπλοκων ορίων απόφασης.
\end{itemize}

\clearpage

% ======================================================
\section{Μέρος 4: Expectation–Maximization (EM)}
% ======================================================

% ------------------------------------------------------
\subsection{Θεωρητικό πλαίσιο: ML με κρυφές μεταβλητές}
% ------------------------------------------------------

Ο στόχος της μέγιστης πιθανοφάνειας (ML) είναι να βρούμε παραμέτρους $\theta$
που μεγιστοποιούν τη log–likelihood των \textbf{παρατηρούμενων} δεδομένων:
\[
\ell(\theta) = \log p(X \mid \theta).
\]

Όταν όμως το μοντέλο περιέχει \textbf{κρυφές (latent) μεταβλητές} $Z$,
γράφουμε
\[
p(X \mid \theta) = \sum_Z p(X,Z \mid \theta),
\]
οπότε
\[
\ell(\theta) = \log \sum_Z p(X,Z \mid \theta).
\]
Η log ενός αθροίσματος είναι δύσκολο να βελτιστοποιηθεί άμεσα,
σε αντίθεση με την περίπτωση \textbf{πλήρων δεδομένων}, όπου θα είχαμε
\[
\log p(X,Z \mid \theta).
\]

Βασικές έννοιες:

\begin{itemize}
  \item \textbf{Complete data}:
        τι θα βλέπαμε αν γνωρίζαμε και τις κρυφές μεταβλητές $Z$.
  \item \textbf{Incomplete / observed data}:
        μόνο τα $X$ (κάποια πεδία λείπουν ή είναι latent).
  \item \textbf{Κρυφές μεταβλητές} $Z$:
        συμπληρώνουν τα δεδομένα έτσι ώστε να απλοποιείται η ML.
\end{itemize}

Η ιδέα του EM είναι να χρησιμοποιήσει την \textbf{αναμενόμενη}
log–likelihood των πλήρων δεδομένων:
\[
Q(\theta,\theta^{\text{old}})
= \mathbb{E}_{Z \mid X,\theta^{\text{old}}}
  \big[ \log p(X,Z \mid \theta) \big].
\]

\begin{description}
  \item[E-step:]
    Υπολογίζουμε τη συνάρτηση
    \[
      Q(\theta,\theta^{\text{old}})
      = \mathbb{E}_{Z \mid X,\theta^{\text{old}}}
        \big[ \log p(X,Z \mid \theta) \big],
    \]
    δηλαδή την αναμενόμενη \emph{πλήρη} log–likelihood,
    όπου η expectation είναι ως προς την τρέχουσα posterior
    των latent $Z$.
  \item[M-step:]
    Επιλέγουμε
    \[
      \theta^{\text{new}} = \arg\max_{\theta} Q(\theta,\theta^{\text{old}}),
    \]
    σαν να είχαμε πραγματικά τα πλήρη δεδομένα.
\end{description}

Επαναλαμβάνουμε τα βήματα E–M μέχρι σύγκλιση. Βαθύτερο θεωρητικό αποτέλεσμα:
κάθε πλήρης επανάληψη EM \textbf{δεν μειώνει} την $\log p(X\mid\theta)$
(μονοτονική αύξηση μέχρι κάποιο τοπικό μέγιστο).

Τυπικές εφαρμογές:
\begin{itemize}
  \item Μείγματα κατανομών (mixture models).
  \item HMM (Baum–Welch).
  \item Προβλήματα με \textbf{ελλιπή δεδομένα} (missing values).
\end{itemize}

% ------------------------------------------------------
\subsection{Το συγκεκριμένο παράδειγμα: εκθετική \texorpdfstring{$\times$}{x} ομοιόμορφη}
% ------------------------------------------------------

Στην άσκηση, τα δεδομένα είναι:
\[
D = \{(1,2),\ (4,5),\ (2,\ast)\},
\]
όπου το $\ast$ δηλώνει άγνωστη τιμή για $x_2$ στο τρίτο δείγμα.

Υποθέτουμε ανεξαρτησία:
\[
p(x_1,x_2 \mid \theta) = p(x_1\mid\theta_1)\,p(x_2\mid\theta_2),
\]
με:

\begin{itemize}
  \item \textbf{Εκθετική} για $x_1$ (με μέσο $\theta_1$):
        \[
        p(x_1 \mid \theta_1) =
        \begin{cases}
          \dfrac{1}{\theta_1} \exp\!\left(-\dfrac{x_1}{\theta_1}\right),
            & x_1 \ge 0,\\[4pt]
          0, & \text{αλλιώς.}
        \end{cases}
        \]

  \item \textbf{Ομοιόμορφη} για $x_2$:
        \[
        p(x_2 \mid \theta_2) = U(0,\theta_2) =
        \begin{cases}
          \dfrac{1}{\theta_2}, & 0 \le x_2 \le \theta_2,\\[4pt]
          0, & \text{αλλιώς,}
        \end{cases}
        \]
        όπου $\theta_2 > 0$.
\end{itemize}

Παράμετροι:
\[
\theta = (\theta_1,\theta_2).
\]
Κρυφή μεταβλητή:
\[
Z := x_2^{(3)},
\]
η άγνωστη τιμή του $x_2$ στο τρίτο δείγμα.

Η άσκηση ζητά να εφαρμόσουμε ένα \textbf{ένα} βήμα EM, με αρχική εκτίμηση
\[
\theta^{(0)} = (\theta_1^{(0)},\theta_2^{(0)}) = (2,3).
\]

% ------------------------------------------------------
\subsection{(α) E-step: κατασκευή του \texorpdfstring{$Q(\theta,\theta^{(0)})$}{Q}}
% ------------------------------------------------------

\paragraph{Πλήρη δεδομένα.}
Αν γνωρίζαμε και το $Z$, τα πλήρη δεδομένα θα ήταν:
\[
\tilde{D} = \{(1,2),\ (4,5),\ (2,Z)\}.
\]

Λόγω ανεξαρτησίας:
\[
p(\tilde{D} \mid \theta)
= \prod_{n=1}^3 p(x_1^{(n)} \mid \theta_1)\, p(x_2^{(n)} \mid \theta_2).
\]
Άρα
\[
\log p(\tilde{D} \mid \theta)
= \sum_{n=1}^3 \log p(x_1^{(n)}\mid\theta_1)
+ \sum_{n=1}^3 \log p(x_2^{(n)}\mid\theta_2).
\]

\paragraph{Μέρος για $x_1$ (χωρίς κρυφές μεταβλητές).}

Τα $x_1$ είναι $1,4,2$. Για την εκθετική:
\[
\log p(x_1^{(n)} \mid \theta_1)
= -\log \theta_1 - \frac{x_1^{(n)}}{\theta_1}.
\]
Άρα:
\[
\sum_{n=1}^3 \log p(x_1^{(n)} \mid \theta_1)
= -3\log\theta_1 - \frac{1+4+2}{\theta_1}
= -3\log\theta_1 - \frac{7}{\theta_1}.
\]
Δεν εξαρτάται από το $Z$, άρα το μέρος αυτό είναι \textbf{ίδιο}
με απλό ML.

\paragraph{Μέρος για $x_2$ (εξαρτάται από $Z$).}

Τα πλήρη $x_2$ θα ήταν $2,5,Z$.
Για την $U(0,\theta_2)$:
\[
\log p(x_2^{(n)} \mid \theta_2)
= -\log\theta_2 + \log \mathbf{1}_{[0,\theta_2]}(x_2^{(n)}),
\]
όπου $\mathbf{1}_{[0,\theta_2]}$ είναι η ενδεικτική του διαστήματος.
Συνεπώς:
\[
\sum_{n=1}^3 \log p(x_2^{(n)} \mid \theta_2)
= -3\log\theta_2 + \sum_{n=1}^3 \log \mathbf{1}_{[0,\theta_2]}(x_2^{(n)}).
\]

Για να μην γίνει η log–likelihood $-\infty$, πρέπει
\[
0\le 2,5,Z \le \theta_2,
\]
δηλαδή
\[
\theta_2 \ge \max\{2,5,Z\}.
\]

\paragraph{Ορισμός του $Q(\theta,\theta^{(0)})$.}

Με βάση τον ορισμό:
\[
Q(\theta,\theta^{(0)})
= \mathbb{E}_{Z \mid D_{\text{obs}},\theta^{(0)}}
  \big[ \log p(\tilde{D}\mid \theta) \big].
\]

Ο παράγοντας για $x_1$ δεν εξαρτάται από $Z$, άρα:
\[
Q(\theta,\theta^{(0)}) = Q_1(\theta_1) + Q_2(\theta_2,\theta_2^{(0)}),
\]
με
\[
Q_1(\theta_1) = -3\log\theta_1 - \frac{7}{\theta_1}.
\]

Για το $Q_2$:
\[
Q_2(\theta_2,\theta_2^{(0)})
= \mathbb{E}_{Z \mid \theta_2^{(0)}}
  \left[ -3\log\theta_2 + \sum_{n=1}^3
         \log \mathbf{1}_{[0,\theta_2]}(x_2^{(n)}) \right].
\]

Η κατανομή του $Z$ υπό $\theta^{(0)}$ είναι:
\[
Z \sim U(0,\theta_2^{(0)}) = U(0,3).
\]

Αν επιλέξουμε $\theta_2 < 5$ ή $\theta_2 < \theta_2^{(0)}$, τότε υπάρχει
μη μηδενική πιθανότητα το $Z$ να βρεθεί εκτός $[0,\theta_2]$, οπότε
η αναμενόμενη log–likelihood γίνεται $-\infty$.
Συνεπώς, για να είναι $Q_2$ πεπερασμένο πρέπει:
\[
\theta_2 \ge \max\{5,\theta_2^{(0)}\}.
\]

Σε αυτό το επιτρεπτό πεδίο, οι ενδεικτικές γίνονται $1$ σχεδόν σίγουρα
ως προς $Z$, άρα:
\[
Q_2(\theta_2,\theta_2^{(0)})
= -3\log\theta_2 + \text{σταθερά ανεξάρτητη από }\theta_2.
\]

Συνολικά (αγνοώντας σταθερές ως προς τη μεγιστοποίηση):
\[
Q(\theta,\theta^{(0)})
= -3\log\theta_1 - \frac{7}{\theta_1}
  - 3\log\theta_2,
\]
με περιορισμό $\theta_2 \ge \max\{5,\theta_2^{(0)}\}$.

% ------------------------------------------------------
\subsection{(β) M-step: ML εκτίμηση παραμέτρων}
% ------------------------------------------------------

\paragraph{Βελτιστοποίηση ως προς $\theta_1$.}

Με
\[
Q_1(\theta_1) = -3\log\theta_1 - \frac{7}{\theta_1},
\]
έχουμε:
\[
\frac{\partial Q_1}{\partial \theta_1}
= -\frac{3}{\theta_1} + \frac{7}{\theta_1^2}.
\]

Θέτοντας ίσο με μηδέν:
\[
-\frac{3}{\theta_1} + \frac{7}{\theta_1^2} = 0
\;\Rightarrow\;
-3\theta_1 + 7 = 0
\;\Rightarrow\;
\boxed{\theta_1^{\text{new}} = \frac{7}{3}.}
\]

Θεωρητικά:
\[
\theta_1^{\text{ML}} = \text{μέση τιμή των }x_1
= \frac{1+4+2}{3} = \frac{7}{3},
\]
δηλαδή ακριβώς ML για εκθετική με παράμετρο κλίμακας.

\paragraph{Βελτιστοποίηση ως προς $\theta_2$.}

Μέσα στο επιτρεπτό πεδίο $\theta_2 \ge \max\{5,\theta_2^{(0)}\}$,
έχουμε:
\[
Q_2(\theta_2) = -3\log\theta_2 + \text{σταθερά}.
\]

Επειδή $-\log\theta_2$ είναι φθίνουσα συνάρτηση, το $Q_2$ μεγιστοποιείται
στη \textbf{μικρότερη} επιτρεπτή τιμή:
\[
\boxed{\theta_2^{\text{new}} = \max\{5,\theta_2^{(0)}\}.}
\]

Με $\theta_2^{(0)} = 3$:
\[
\boxed{\theta_2^{(1)} = 5.}
\]

Συνολική ενημέρωση:
\[
\theta^{(1)} = \left(\frac{7}{3},\, 5\right).
\]

Παρατηρούμε:

\begin{itemize}
  \item Για την εκθετική, η EM απλώς \textbf{συμπίπτει} με το κλασικό ML.
  \item Για την ομοιόμορφη, η EM ``μαθαίνει'' ότι η στήριξη πρέπει
        να καλύπτει τη μέγιστη παρατήρηση $x_2=5$, άρα
        $\theta_2^{\text{ML}} = \max x_2^{(n)} = 5$.
\end{itemize}

% ------------------------------------------------------
\subsection{(γ) Γραφική ερμηνεία και σύνδεση με άλλα μοντέλα}
% ------------------------------------------------------

\paragraph{Πριν το EM: $\theta^{(0)}=(2,3)$.}

\begin{itemize}
  \item $p(x_1\mid \theta_1^{(0)}) = \frac{1}{2} e^{-x_1/2}$:
        εκθετική με μέσο $2$ (πιο ``απότομη'').
  \item $p(x_2\mid \theta_2^{(0)}) = U(0,3)$:
        \[
        p(x_2) =
        \begin{cases}
          \dfrac{1}{3}, & 0\le x_2\le 3,\\[4pt]
          0, & \text{αλλιώς.}
        \end{cases}
        \]
        Δεν καλύπτει τη τιμή $x_2=5$.
\end{itemize}

\paragraph{Μετά το EM: $\theta^{(1)}=(7/3,5)$.}

\begin{itemize}
  \item $p(x_1\mid \theta_1^{(1)}) = \frac{3}{7} e^{-3x_1/7}$:
        μέσος $\frac{7}{3} \approx 2.33$, άρα η κατανομή ``απλώνεται'' ελαφρώς.
  \item $p(x_2\mid \theta_2^{(1)}) = U(0,5)$:
        \[
        p(x_2) =
        \begin{cases}
          \dfrac{1}{5}, & 0\le x_2\le 5,\\[4pt]
          0, & \text{αλλιώς,}
        \end{cases}
        \]
        το διάστημα μεγαλώνει ώστε να περιλάβει τη μέγιστη παρατήρηση,
        και το ύψος της πυκνότητας μειώνεται αντίστοιχα.
\end{itemize}

\paragraph{Σύνδεση με άλλα μέρη.}

\begin{itemize}
  \item Όπως στο Μέρος~1 κάνουμε ML για Γκαουσιανές,
        εδώ κάνουμε ML με \textbf{latent} ποσότητες, άρα χρειάζεται EM.
  \item Όπως στα HMM (Μέρος~2) τα κρυφά states $q_t$ οδηγούν σε forward–backward,
        εδώ το κρυφό $Z$ οδηγεί σε E-step με expectation ως προς $p(Z\mid X,\theta^{(0)})$.
  \item Η φιλοσοφία ``latent + EM'' είναι κοινή για \textbf{μείγματα},
        \textbf{HMM} και μοντέλα με ελλιπή δεδομένα.
\end{itemize}

\clearpage

% ======================================================
\section{Μέρος 5: Bayes Network – Μοντέλο Chest Clinic}
% ======================================================

% ------------------------------------------------------
\subsection{Θεωρία Bayesian Networks}
% ------------------------------------------------------

Ένα \textbf{Bayesian Network (BN)} είναι ένας κατευθυνόμενος ακυκλικός γράφος (DAG)
που περιγράφει την joint κατανομή ενός συνόλου τυχαίων μεταβλητών
$\{X_1,\dots,X_n\}$ μέσω:
\[
P(X_1,\dots,X_n) = \prod_{i=1}^n P(X_i \mid \text{γονείς}(X_i)).
\]

Κύρια χαρακτηριστικά:

\begin{itemize}
  \item Κάθε ακμή $U\to V$ δηλώνει \emph{άμεση πιθανοτική εξάρτηση}.
  \item Η δομή του γράφου κωδικοποιεί \textbf{συνθήκες ανεξαρτησίας}
        μέσω της έννοιας \textbf{d-separation}.
  \item Το BN επιτρέπει:
        \begin{itemize}
          \item \textbf{forward} reasoning (από αιτίες σε συμπτώματα),
          \item \textbf{diagnostic} reasoning (από συμπτώματα σε αιτίες),
          \item \textbf{intercausal} επεξηγήσεις (π.χ. colliders).
        \end{itemize}
\end{itemize}

Τρεις βασικοί τύποι μικρο-δομών:

\begin{itemize}
  \item \emph{Chain:} $X \to Y \to Z$.
  \item \emph{Fork:} $X \leftarrow Y \to Z$.
  \item \emph{Collider:} $X \to Y \leftarrow Z$.
\end{itemize}

Η θεωρία d-separation λέει πότε δύο κόμβοι είναι συνθηκά ανεξάρτητοι
δεδομένου ενός συνόλου $Z$:

\begin{itemize}
  \item Αν σε κάθε path ανάμεσά τους υπάρχει:
        \begin{itemize}
          \item μη-collider κόμβος που ανήκει στο $Z$, ή
          \item collider κόμβος τέτοιος ώστε ούτε αυτός ούτε απόγονός του
                ανήκει στο $Z$,
        \end{itemize}
        τότε οι κόμβοι είναι \textbf{d-separated} και άρα συνθηκά ανεξάρτητοι.
\end{itemize}

% ------------------------------------------------------
\subsection{Το μοντέλο Chest Clinic ως παράδειγμα BN}
% ------------------------------------------------------

Το Chest Clinic BN μοντελοποιεί απλοποιημένη ιατρική διάγνωση
για αναπνευστικά προβλήματα.

Κόμβοι:

\begin{itemize}
  \item $a$: επίσκεψη στην περιοχή Ω (visit to Asia).
  \item $s$: καπνιστής (smoker).
  \item $t$: φυματίωση (tuberculosis).
  \item $l$: καρκίνος του πνεύμονα (lung cancer).
  \item $b$: βρογχίτιδα (bronchitis).
  \item $e$: ``either'' (τουλάχιστον ένα από $t,l$ είναι true).
  \item $x$: αποτέλεσμα ακτινογραφίας (X-ray).
  \item $d$: δύσπνοια (dyspnoea).
\end{itemize}

Δομή:

\begin{itemize}
  \item $a \to t$ (visit επηρεάζει την πιθανότητα φυματίωσης).
  \item $s \to l$, $s \to b$ (smoking επηρεάζει cancer \& bronchitis).
  \item $t \to e$, $l \to e$ (either αν υπάρχει t, l ή και τα δύο).
  \item $e \to x$, $e \to d$ (η ύπαρξη σοβαρής νόσου επηρεάζει ακτινογραφία και δύσπνοια).
  \item $b \to d$ (bronchitis επηρεάζει δύσπνοια).
\end{itemize}

% ------------------------------------------------------
\subsection{(α) Παραγοντοποίηση της joint κατανομής}
% ------------------------------------------------------

Από τον ορισμό BN:
\[
P(a,s,t,l,b,e,x,d)
= P(a)\,P(s)\,P(t \mid a)\,P(l \mid s)\,P(b \mid s)\,
  P(e \mid t,l)\,P(x \mid e)\,P(d \mid e,b).
\]

Αυτό είναι το \textbf{θεωρητικό πρότυπο}:
κάθε κόμβος συνοδεύεται από έναν παράγοντα $P(\text{κόμβος}\mid\text{γονείς})$.

% ------------------------------------------------------
\subsection{(β) Συνθηκές ανεξαρτησίας και d-separation}
% ------------------------------------------------------

Θεωρούμε δύο δηλώσεις:

\begin{enumerate}
  \item $l \perp\!\!\!\perp b \mid s$;
  \item $a \perp\!\!\!\perp s \mid l$;
\end{enumerate}

\paragraph{1) $l \perp\!\!\!\perp b \mid s$;}

Ο βασικός path είναι:
\[
l \leftarrow s \rightarrow b,
\]
δηλαδή ένας \textbf{fork} με ``γιαγιά'' τον $s$.

\begin{itemize}
  \item Χωρίς conditioning: $l$ και $b$ είναι εξαρτημένα μέσω του $s$
        (το smoking επηρεάζει και τα δύο).
  \item Με conditioning στο $s$: σε fork, αν \textbf{γνωρίζουμε}
        την κοινή αιτία, οι απόγονοι γίνονται συνθηκά ανεξάρτητοι.
\end{itemize}

Άρα:
\[
\boxed{l \perp\!\!\!\perp b \mid s.}
\]

\paragraph{2) $a \perp\!\!\!\perp s \mid l$;}

Κοιτάμε paths από $a$ σε $s$.
Ένας σημαντικός είναι:
\[
a \to t \to e \leftarrow l \leftarrow s.
\]

Υπάρχει collider στη θέση $e$ ($t\to e \leftarrow l$).
Ο $l$ είναι μη-collider στο path (βέλος από $s$ και προς $e$).

\begin{itemize}
  \item Χωρίς conditioning: ο collider $e$ \textbf{μπλοκάρει} το path,
        άρα $a$ και $s$ είναι ανεξάρτητοι.
  \item Με conditioning στο $l$ (μη-collider) εξακολουθούμε να έχουμε
        μπλοκάρισμα στον $e$ (δεν παρατηρούμε τον collider ούτε απόγονό του),
        άρα το path παραμένει μπλοκαρισμένο.
\end{itemize}

Δεν υπάρχουν άλλοι ενεργοί paths που να ``ανοίγουν'' εξάρτηση,
άρα:
\[
\boxed{a \perp\!\!\!\perp s \mid l.}
\]

Ερμηνεία: αν ήδη γνωρίζουμε αν ο ασθενής έχει καρκίνο ($l$),
η πληροφορία για το αν επισκέφθηκε την περιοχή Ω ($a$) δεν μας δίνει
παραπάνω πληροφορία για το αν είναι καπνιστής ($s$), και αντίστροφα.

% ------------------------------------------------------
\subsection{(γ) Υπολογισμός πιθανοτήτων με product–sum}
% ------------------------------------------------------

Η factorization
\[
P(a,s,t,l,b,e,x,d)
= P(a)\,P(s)\,P(t \mid a)\,P(l \mid s)\,P(b \mid s)\,
  P(e \mid t,l)\,P(x \mid e)\,P(d \mid e,b)
\]
μας επιτρέπει να υπολογίσουμε οποιαδήποτε περιθωριακή ή συνθήκη
πιθανότητα με τον γενικό κανόνα \textbf{sum–product}:

\[
P(\text{μεταβλητές ενδιαφέροντος})
= \sum_{\text{λοιπές μεταβλητές}} \prod_{\text{παράγοντες } \phi}
  \phi(\text{μεταβλητές } \phi).
\]

Ενδεικτικά:

\begin{itemize}
  \item
    \[
    P(x=T)
    = \sum_{a,s,t,l,b,e,d}
      P(a)\,P(s)\,P(t \mid a)\,P(l \mid s)\,P(b \mid s)\,
      P(e \mid t,l)\,P(x=T \mid e)\,P(d \mid e,b).
    \]
  \item
    \[
    P(x=T\mid t=F)
    = \frac{P(x=T,t=F)}{P(t=F)},
    \]
    όπου
    \[
    P(x=T,t=F)
    = \sum_{a,s,l,b,e,d}
      P(a,s,t=F,l,b,e,x=T,d).
    \]
  \item
    \[
    P(t=T\mid x=F)
    = \frac{P(t=T,x=F)}{P(x=F)}.
    \]
  \item
    \[
    P(t=T\mid x=F,d=T)
    = \frac{P(t=T,x=F,d=T)}{P(x=F,d=T)}.
    \]
\end{itemize}

Στην πράξη, οι αριθμητικοί υπολογισμοί γίνονται με
\textbf{message passing} σε κατάλληλη δομή (junction tree), όχι με
«ωμό» άθροισμα σε όλες τις $2^8$ περιπτώσεις.

% ------------------------------------------------------
\subsection{Moralization, triangulation και αποδοτικό inference}
% ------------------------------------------------------

Για μεγάλα δίκτυα, το απλό sum–product είναι εκθετικό. Η κλασική
θεωρητική διαδικασία για αποδοτικό exact inference είναι:

\begin{enumerate}
  \item \textbf{Moralization}:
        \begin{itemize}
          \item Ενώνουμε (παντρεύουμε) όλους τους γονείς κάθε κόμβου
                (π.χ.\ $t$–$l$ για τον $e$).
          \item Κάνουμε τον γράφο μη κατευθυνόμενο.
        \end{itemize}
  \item \textbf{Triangulation}:
        \begin{itemize}
          \item Προσθέτουμε ακμές ώστε να μην υπάρχουν κύκλοι $\ge 4$
                χωρίς χορδή.
          \item Προκύπτουν \textbf{cliques}, στις οποίες ορίζουμε potentials.
        \end{itemize}
  \item \textbf{Junction tree \& message passing}:
        \begin{itemize}
          \item Χτίζουμε ένα δέντρο από cliques.
          \item Πολλαπλασιάζουμε τους παράγοντες $P(\cdot\mid\cdot)$
                στις κατάλληλες cliques.
          \item Εφαρμόζουμε sum–product (messages) για να πάρουμε όλες τις
                ζητούμενες περιθωριακές/posterior πιθανότητες.
        \end{itemize}
\end{enumerate}

Η λογική είναι ανάλογη με forward–backward στα HMM:
\textbf{τοπικά προϊόντα + άθροισμα πάνω σε κρυφές μεταβλητές}.

% ------------------------------------------------------
\subsection{Σύνδεση με τα υπόλοιπα μέρη}
% ------------------------------------------------------

\begin{itemize}
  \item Όπως στον Bayes ταξινομητή (Μέρος~1) χρειαζόμαστε $P(C\mid x)$,
        εδώ χρειαζόμαστε γενικές $P(\text{μέρος του γράφου} \mid \text{evidence})$.
  \item Όπως στο EM (Μέρος~4) γράφουμε joint κατανομές ως γινόμενο
        απλών παραγόντων και αθροίζουμε πάνω σε latent μεταβλητές,
        στα BN κάνουμε sum–product πάνω σε παράγοντες $P(X_i\mid \text{γονείς})$.
  \item Η ιδέα \textbf{τοπικών πιθανοτήτων + δομή γραφήματος}
        είναι κοινή σε HMM, BN, factor graphs και γενικά σε γραφικά μοντέλα.
\end{itemize}

\clearpage

% ======================================================
\section{Μέρος 6: Απλό RNN που Μετράει Μονάδες}
% ======================================================

% ------------------------------------------------------
\subsection{RNN ως δυναμικό διακριτικό μοντέλο}
% ------------------------------------------------------

Ένα \textbf{Recurrent Neural Network (RNN)} είναι ένα νευρωνικό δίκτυο
που επεξεργάζεται \textbf{ακολουθίες} εισόδων
\[
x_1,x_2,\dots,x_T
\]
και διατηρεί μια \textbf{εσωτερική κατάσταση} $s_k$ στο χρόνο $k$.

Γενική ιδέα:
\[
s_k = f(s_{k-1},x_k;\theta),
\quad
y_k = g(s_k;\theta),
\]
όπου $f,g$ είναι συνήθως μη γραμμικές συναρτήσεις (π.χ.\ tanh, ReLU κ.λπ.).

\begin{itemize}
  \item Το RNN είναι \textbf{διακριτικό} μοντέλο:
        μαθαίνει απευθείας $y$ (ή $P(y\mid x_{1:T})$), όχι $p(x_{1:T})$.
  \item Σε αντίθεση με τα HMM (γενετικά πιθανοτικά μοντέλα), τα RNN
        έχουν \emph{ντετερμινιστική} κατάσταση με βάρη.
  \item Σε αντίθεση με MLP (feedforward), τα RNN είναι \textbf{δυναμικά}:
        η ίδια συνάρτηση $f$ εφαρμόζεται επαναληπτικά στο χρόνο.
\end{itemize}

% ------------------------------------------------------
\subsection{Ειδικό μοντέλο: γραμμικό RNN που μετράει μονάδες}
% ------------------------------------------------------

Στην άσκηση, έχουμε μια ιδιαίτερα απλή μορφή RNN:

\[
s_k = w_{\mathrm{rec}}\, s_{k-1} + w_x\, x_k,\quad s_0=0,
\]
όπου:
\begin{itemize}
  \item $x_k \in \{0,1\}$ – δυαδική είσοδος,
  \item $s_k$ – scalar hidden state,
  \item $w_{\mathrm{rec}}$ – recurrent βάρος,
  \item $w_x$ – input βάρος.
\end{itemize}

Η έξοδος ορίζεται ως:
\[
y = s_T,
\]
δηλαδή μόνο στο τελευταίο timestep.

Ο στόχος είναι:
\[
t = \sum_{k=1}^T x_k,
\]
δηλαδή το \textbf{πλήθος μονάδων} στην ακολουθία.

Κόστος:
\[
L(y,t) = (t - y)^2 = (t - s_T)^2.
\]

Θεωρητικά, αν πάρουμε $w_{\mathrm{rec}}=1$, $w_x=1$, τότε
\[
s_k = s_{k-1} + x_k \Rightarrow s_k = \sum_{i=1}^k x_i,
\]
άρα $y = s_T = t$: το RNN γίνεται ιδανικός \emph{accumulator} μονάδων.

% ------------------------------------------------------
\subsection{(α) Backpropagation Through Time (BPTT) και gradients}
% ------------------------------------------------------

Θέλουμε τις παραγώγους:
\[
\frac{\partial L}{\partial w_{\mathrm{rec}}},\quad
\frac{\partial L}{\partial w_x},
\]
για να εφαρμόσουμε gradient descent:
\[
w_{\mathrm{rec}} \leftarrow w_{\mathrm{rec}} - \eta \frac{\partial L}{\partial w_{\mathrm{rec}}},\quad
w_x \leftarrow w_x - \eta \frac{\partial L}{\partial w_x}.
\]

Ο αλγόριθμος είναι \textbf{Backpropagation Through Time} (BPTT):
ξετυλίγουμε το RNN στον χρόνο και εφαρμόζουμε backprop όπως σε πολύ
βαθύ δίκτυο με κοινά βάρη.

\paragraph{1. Παράγωγος κόστους ως προς $s_T$.}

Με $L = (t - s_T)^2$:
\[
\delta_T := \frac{\partial L}{\partial s_T} = 2(s_T - t).
\]

\paragraph{2. Παράγωγος προς προηγούμενες καταστάσεις.}

Για $k<T$, το $L$ εξαρτάται από $s_k$ μέσω του $s_{k+1}$:
\[
s_{k+1} = w_{\mathrm{rec}} s_k + w_x x_{k+1}.
\]

Άρα:
\[
\delta_k := \frac{\partial L}{\partial s_k}
= \frac{\partial L}{\partial s_{k+1}} \frac{\partial s_{k+1}}{\partial s_k}
= \delta_{k+1} \, w_{\mathrm{rec}}.
\]

Επαναληπτικά:
\[
\delta_k = \delta_T\, w_{\mathrm{rec}}^{\,T-k}.
\]

\paragraph{3. Gradient ως προς $w_{\mathrm{rec}}$.}

Σε κάθε timestep $k$, η εξάρτηση από $w_{\mathrm{rec}}$ είναι
\[
s_k = w_{\mathrm{rec}} s_{k-1} + w_x x_k
\Rightarrow
\frac{\partial s_k}{\partial w_{\mathrm{rec}}} = s_{k-1}.
\]

Συνολικά:
\[
\boxed{
\frac{\partial L}{\partial w_{\mathrm{rec}}}
= \sum_{k=1}^T \delta_k\, s_{k-1}.
}
\]

\paragraph{4. Gradient ως προς $w_x$.}

Παρομοίως:
\[
\frac{\partial s_k}{\partial w_x} = x_k,
\]
άρα:
\[
\boxed{
\frac{\partial L}{\partial w_x}
= \sum_{k=1}^T \delta_k\, x_k.
}
\]

\paragraph{5. Κανόνες ανανέωσης.}

Με gradient descent:
\[
\boxed{
\begin{aligned}
w_{\mathrm{rec}} &\leftarrow
    w_{\mathrm{rec}} - \eta \sum_{k=1}^T \delta_k\, s_{k-1},\\[4pt]
w_x &\leftarrow
    w_x - \eta \sum_{k=1}^T \delta_k\, x_k,
\end{aligned}
}
\]
όπου
\[
\delta_T = 2(s_T - t),\quad
\delta_k = w_{\mathrm{rec}}\, \delta_{k+1},\ k=T-1,\dots,1.
\]

Αυτό είναι ένα πλήρες παράδειγμα BPTT για ένα γραμμικό RNN.

% ------------------------------------------------------
\subsection{(β) Παράδειγμα ενός update με αρχικά βάρη}
% ------------------------------------------------------

Για την ακολουθία
\[
[x_1,\dots,x_6] = [1,0,1,0,1,0],\quad T=6,\quad t=3,
\]
και αρχικά βάρη
\[
w_{\mathrm{rec}}^{(0)}=1,\quad w_x^{(0)}=0,
\]
ισχύει:

\begin{itemize}
  \item \textbf{Forward:}
        \[
        s_k^{(0)} = 1\cdot s_{k-1}^{(0)} + 0\cdot x_k = 0
        \quad\Rightarrow\quad s_k^{(0)}=0,\;\forall k.
        \]
        Άρα $y^{(0)}=0$, $L^{(0)}=(3-0)^2=9$.
  \item \textbf{Backward:}
        \[
        \delta_6^{(0)} = 2(s_6^{(0)}-t) = -6,\quad
        \delta_k^{(0)} = w_{\mathrm{rec}}^{(0)} \delta_{k+1}^{(0)}
        = -6,\;\forall k.
        \]
  \item \textbf{Gradients:}
        \[
        \frac{\partial L}{\partial w_x}\Big|_{(0)}
        = \sum_{k=1}^6 \delta_k^{(0)} x_k
        = -6\cdot(1+0+1+0+1+0) = -18,
        \]
        \[
        \frac{\partial L}{\partial w_{\mathrm{rec}}}\Big|_{(0)}
        = \sum_{k=1}^6 \delta_k^{(0)} s_{k-1}^{(0)} = 0,
        \]
        επειδή όλα τα $s_{k-1}^{(0)}=0$.
  \item \textbf{Update:}
        \[
        w_x^{(1)} = 0 - \eta (-18) = 18\eta,\quad
        w_{\mathrm{rec}}^{(1)} = 1.
        \]
\end{itemize}

Έτσι, το πρώτο update ``ανοίγει'' τον δρόμο ώστε το δίκτυο να αρχίσει
να μετράει τις μονάδες μέσω του $w_x$.

Τα επόμενα updates υπολογίζονται με την ίδια διαδικασία, τώρα όμως
τα $s_k$ είναι μη μηδενικά και το gradient ως προς $w_{\mathrm{rec}}$
γίνεται επίσης μη μηδενικό.

% ------------------------------------------------------
\subsection{(γ) Καλύτερη αρχικοποίηση και ιδανική λύση}
% ------------------------------------------------------

\paragraph{Ιδανικά βάρη.}
Για τέλειο counting:
\[
w_{\mathrm{rec}}^\ast = 1,\quad
w_x^\ast = 1.
\]
Τότε:
\[
s_k^\ast = \sum_{i=1}^k x_i,\quad
y^\ast = s_T^\ast = t,
\]
οπότε $L=0$ και όλα τα gradients γίνονται μηδενικά.

\paragraph{Καλή αρχικοποίηση στην πράξη.}
Συνήθως δεν γνωρίζουμε τη βέλτιστη λύση, αλλά:
\[
w_{\mathrm{rec}}^{(0)} \approx 1,\quad w_x^{(0)} \approx 1
\]
είναι λογική επιλογή για \emph{accumulator} τύπου RNN:

\begin{itemize}
  \item $|w_{\mathrm{rec}}|<1$ οδηγεί σε \textbf{vanishing} μνήμη και gradients,
        ειδικά για μεγάλα $T$.
  \item $|w_{\mathrm{rec}}|>1$ οδηγεί σε \textbf{exploding} καταστάσεις και gradients.
  \item $w_{\mathrm{rec}}\approx 1$ διατηρεί τις πληροφορίες για αρκετά βήματα,
        χωρίς να εκρήγνυνται οι τιμές.
\end{itemize}

% ------------------------------------------------------
\subsection{(δ) Vanishing / exploding gradients και συγκρίσεις}
% ------------------------------------------------------

Η σχέση
\[
\delta_k = w_{\mathrm{rec}}^{\,T-k}\, \delta_T
\]
δείχνει καθαρά:

\begin{itemize}
  \item Αν $|w_{\mathrm{rec}}|<1$, τότε $|w_{\mathrm{rec}}^{\,T-k}|$ μικραίνει
        εκθετικά με το μήκος της ακολουθίας: \textbf{vanishing gradients}.
  \item Αν $|w_{\mathrm{rec}}|>1$, τότε μεγαλώνει εκθετικά:
        \textbf{exploding gradients}.
\end{itemize}

Γι’ αυτό τα πιο εξελιγμένα RNN (LSTM, GRU) χρησιμοποιούν gates,
skip connections και άλλες τεχνικές για να αντιμετωπίσουν αυτά τα φαινόμενα.

Συγκριτικά:

\begin{itemize}
  \item Σε σχέση με \textbf{LMS / Perceptron} (Μέρος~1),
        εδώ η δομή είναι δυναμική στον χρόνο, αλλά η βασική ιδέα
        gradient descent είναι η ίδια.
  \item Σε σχέση με \textbf{HMM} (Μέρος~2),
        τα RNN δεν έχουν explicit κατανομές μετάβασης/εκπομπής,
        αλλά μια ντετερμινιστική δυναμική $s_k=f(s_{k-1},x_k)$,
        την οποία μαθαίνουν από τα δεδομένα.
  \item Σε σχέση με \textbf{MLP και backprop} (Μέρος~3),
        το BPTT είναι απλώς backprop σε ένα δίκτυο που έχει επαναληπθεί
        $T$ φορές με κοινά weights.
\end{itemize}

Έτσι, το απλό RNN της άσκησης είναι ένα ιδανικό θεωρητικό παράδειγμα
για να κατανοήσεις:
\begin{itemize}
  \item πώς η πληροφορία ρέει στον χρόνο,
  \item πώς τα gradients πολλαπλασιάζονται ξανά και ξανά με το $w_{\mathrm{rec}}$,
  \item πώς συνδέεται η δυναμική του RNN με τα φαινόμενα vanishing/exploding.
\end{itemize}


\end{document}
